\documentclass{tietreport}

\usepackage{graphicx}
\usepackage{array}
\usepackage{longtable}

\usepackage[
  date=year,
  style=alphabetic,
  backend=biber,
  sorting=nty,
  maxnames=5,
  minnames=3,
  maxalphanames=4,
  minalphanames=3,
  backref=true,
  doi=false,
  isbn=false,
  url=false,
  eprint=false
]{biblatex}
\addbibresource{references.bib}

%% ==================================================
%% TITLE PAGE FIELDS
%% ==================================================

\TitlePageHeader{%
  UCS503 -- Software Engineering%
}

\ProjectTitle{%
  MeshChat%
}

\ProjectSubTitle{%
  Radio Hardware, Integration and Deployment --- Project Journal%
}

\author{%
  \texttt{1024031043} Sanket Kamboj%
}

\TitlePageSubText{%
  Computer Engineering (COE), BE Third Year -- 5th Semester%
}

\AdvisorName{%
  Ms. Jhony Bansal%
}

\TitlePageFooterText{{%
    \large End-Semester Evaluation \\[0.2cm]
    \bfseries Computer Science and Engineering
    Department%
  } \\[0.15cm]
  Thapar Institute of Engineering and Technology,
  Patiala%
}

%% ==================================================

\begin{document}

\maketitle

\tableofcontents

% ===================================================================
\chapter{Scope of This Journal}
% ===================================================================

MeshChat is an infrastructure-free messaging system: phones join the
WiFi network of a nearby ESP32-C3 node, nodes relay messages over a
LoRa radio link to a Raspberry Pi master, and the master routes them
to the node where the recipient is attached. The system is a
single-hop star, not a multi-hop mesh.

This journal is the second of two project journals. My teammate
Parth Deshmukh's journal (folder \texttt{parth\_journal}) follows the
software layers: architecture, the packet protocol, the master's
routing, polling and queue code, the node firmware structure, the
host-side test suites and the integration review. To avoid
repeating it, mine covers what that journal leaves out: the
\textbf{physical hardware}, the \textbf{radio interface on both
platforms}, the \textbf{over-the-air integration} of node and master,
the \textbf{phone interface and its build pipeline}, and
\textbf{deployment and field validation}. Table~\ref{tab:split} shows
the split; the few overlaps are deliberate and are treated from a
different angle.

\begin{table}[h]
\centering
\small
\begin{tabular}{|p{1.7in}|p{1.9in}|p{1.9in}|}
  \hline
  \textbf{Topic} & \textbf{Parth's journal} & \textbf{This journal} \\
  \hline
  Architecture and scope & Star topology, state model &
  Hardware selection, power, RF numbers \\
  \hline
  Radio band mismatch & Noted as a problem &
  Chip limits, frequency words, migration path \\
  \hline
  Packet protocol & Types, fields, validation &
  Byte parity with node and timing over the air \\
  \hline
  Master software & Routing, poller, SQLite queue &
  Linux radio interface (spidev, libgpiod), wiring \\
  \hline
  Node firmware & Tasks, queues, session table &
  SPI bring-up, DIO0 interrupt, WiFi and portal tuning \\
  \hline
  Phone interface & One line (embedded client) &
  Full web client, protocol and build pipeline \\
  \hline
  Testing & Host suites, 43 and 149 results &
  Field-validation plan, what tests cannot show \\
  \hline
  Review & Remaining engineering work &
  Deployment gaps, security findings, fix list \\
  \hline
\end{tabular}
\caption{Division of coverage between the two journals}
\label{tab:split}
\end{table}

All facts here come from the repository: the source code, its
comments, the design document and the database sketch
(\texttt{db.txt}, committed 3 September). Figures marked
\emph{computed} are calculations; the repository holds no field
measurements yet, and this journal says so wherever it matters.

% ===================================================================
\chapter{Hardware and RF Reference}
% ===================================================================

\section{Components}

\begin{itemize}
  \item \textbf{Node:} ESP32-C3 mini board and an SX1278 (Ra-02,
    AI-Thinker) LoRa module, 433~MHz antenna.
  \item \textbf{Master:} Raspberry Pi with a second Ra-02 module on
    SPI0, 433~MHz antenna. The design document recommends a
    separate 3.3~V regulator for the radio, because transmit current
    spikes of about 120~mA exceed what the Pi's 3.3~V header pin is
    sized for.
\end{itemize}

\section{Wiring}

\begin{table}[h]
\centering
\small
\begin{tabular}{|l|l|l|}
  \hline
  \textbf{ESP32-C3} & \textbf{Ra-02} & \textbf{Signal} \\
  \hline
  GPIO 6 & SCK & SPI clock \\
  \hline
  GPIO 7 & MOSI & Data to radio \\
  \hline
  GPIO 2 & MISO & Data from radio \\
  \hline
  GPIO 10 & NSS & Chip select \\
  \hline
  GPIO 3 & RESET & Radio reset \\
  \hline
  GPIO 1 & DIO0 & RxDone / TxDone interrupt \\
  \hline
  3.3 V, GND & 3.3 V, GND & Power (never 5 V) \\
  \hline
\end{tabular}
\caption{Node wiring (all pin numbers live in \texttt{hardware\_config.h})}
\end{table}

\begin{table}[h]
\centering
\small
\begin{tabular}{|l|l|l|l|}
  \hline
  \textbf{Ra-02} & \textbf{BCM GPIO} & \textbf{Pin} & \textbf{Note} \\
  \hline
  SCK & GPIO 11 & 23 & SPI0 clock \\
  \hline
  MOSI & GPIO 10 & 19 & SPI0 \\
  \hline
  MISO & GPIO 9 & 21 & SPI0 \\
  \hline
  NSS & GPIO 22 & 15 & Driven manually \\
  \hline
  RESET & GPIO 25 & 22 & Low 10 ms at start \\
  \hline
  DIO0 & GPIO 23 & 16 & Interrupt line \\
  \hline
  3.3 V, GND & -- & 1 or 17, 6 & Separate regulator advised \\
  \hline
\end{tabular}
\caption{Master wiring; the working pin for DIO0 is GPIO~23 although
\texttt{radio.h} says 24}
\end{table}

\section{Radio Configuration}

Both radios use the same settings, written directly to SX1278
registers \cite{sx1278datasheet}.

\begin{table}[h]
\centering
\small
\begin{tabular}{|l|l|}
  \hline
  \textbf{Setting} & \textbf{Value} \\
  \hline
  Frequency word & \texttt{0x6C4000} (433 MHz) \\
  \hline
  Modem config 1, 2, 3 & \texttt{0x72}, \texttt{0x74}, \texttt{0x04} \\
  & 125 kHz, CR 4/5, SF7, CRC on, AGC auto \\
  \hline
  Preamble, sync word & 8 symbols, \texttt{0x12} (private) \\
  \hline
  FIFO base TX, RX & \texttt{0x80}, \texttt{0x00} \\
  \hline
  PA configuration & Node \texttt{0x8C} (+14 dBm); master \texttt{0x80} (+2 dBm) \\
  \hline
\end{tabular}
\caption{Radio register settings}
\end{table}

The frequency word is $f(\mathrm{MHz}) \times 16384$. This gives
\texttt{0x6C4000} for 433~MHz and \texttt{0xD84400} for 865.0625~MHz,
the first channel of the design document's band.

\section{Airtime and Link Budget (computed)}

Using the time-on-air formula at SF7, 125~kHz, CR 4/5 (symbol time
1.024~ms, preamble 12.5~ms):

\begin{table}[h]
\centering
\small
\begin{tabular}{|l|r|r|r|}
  \hline
  \textbf{Frame} & \textbf{Bytes} & \textbf{Symbols} &
  \textbf{Airtime (ms)} \\
  \hline
  ACK or empty POLL & 4 & 18 & 31.0 \\
  \hline
  PRESENCE & 23 & 48 & 61.7 \\
  \hline
  DATA, 20 characters & 58 & 98 & 112.9 \\
  \hline
  DATA, maximum & 138 & 208 & 225.5 \\
  \hline
\end{tabular}
\caption{Time on air}
\end{table}

A polling slot therefore costs about 62~ms idle, 175~ms for an uplink
message and 144~ms for a downlink delivery; an idle cycle over $N$
nodes takes about $62N$~ms.

\begin{table}[h]
\centering
\small
\begin{tabular}{|l|r|r|r|r|r|}
  \hline
  \textbf{Case} & \textbf{TX} & \textbf{Loss} & \textbf{RX} &
  \textbf{Sens.} & \textbf{Margin} \\
  \hline
  Planned & 20 & 103 & $-83$ & $-123$ & 40 dB \\
  \hline
  Uplink (node) & 14 & 103 & $-89$ & $-123$ & 34 dB \\
  \hline
  Downlink (master) & 2 & 103 & $-101$ & $-123$ & 22 dB \\
  \hline
\end{tabular}
\caption{Link budget at 1~km (dBm, dB; antenna gains ignored)}
\end{table}

\section{The Band Mismatch}

The design document specifies an SX1276 at 865--867~MHz and says the
SX1278 cannot reach that band. The modules in use are SX1278
(about 420--525~MHz), so the prototype runs at 433~MHz. All four chips
of the SX1276--79 family report version \texttt{0x12}, so the
firmware's version check cannot reveal which variant is fitted. The
design document also says 433~MHz belongs to a different regulatory
exemption class, while the hardware configuration comments call it
permitted; this is unresolved and needs a definite answer before
deployment. Migrating to an SX1276 means a new frequency word and
clearing the low-frequency-mode bit.

% ===================================================================
\chapter{Weekly Journal}
% ===================================================================

The weeks follow the same six-week structure as the companion
journal; the full text of each week is in the \texttt{sanket-journal}
folder as \texttt{week\_1.md} to \texttt{week\_6.md}. Here each week is
summarised.

\section{Week 1 -- Hardware Selection, Power Budget and RF Planning}

\paragraph{Done.} Chose the hardware, fixed the RF parameter set,
computed airtime and the link budget, and wrote the first version of
the SQLite queue table (\texttt{db.txt}: destination and source user
ID, message ID, TTL, payload and length, expiry and creation time;
the final \texttt{queue.c} adds a destination index and WAL
journaling).

\paragraph{Problems.} The designed radio band is out of reach for the
modules in hand; the regulatory status of 433~MHz is unresolved;
transmit power on the two sides differs, leaving the downlink 18~dB
weaker than planned.

\paragraph{Outcome.} The deviation is documented with a migration
path, and design claims became numbers (airtime, margin, current)
that later separated software problems from physics.

\section{Week 2 -- Node Radio Wiring and SX1278 Bring-Up}

\paragraph{Done.} Wired the radio to the ESP32-C3; configured SPI2
at 8~MHz, mode 0; implemented the reset pulse, version check
(\texttt{0x12}), register setup and the DIO0 interrupt, which only
notifies the radio task. The radio task alone owns the chip and feeds
the watchdog in two-second slices while waiting for a POLL.

\paragraph{Problems.} The LoRa bit is writable only in sleep mode; a
floating DIO0 caused edge storms; TxDone and RxDone share one pin, so
the mapping must follow the current mode. The logged RSSI uses an
offset of $-157$, which belongs to the high-frequency port; at 433~MHz
the offset is $-164$, so logged RSSI is about 7~dB optimistic.

\paragraph{Outcome.} A fixed sleep--LoRa--standby sequence, a
pull-down and a yield on empty interrupts. Lesson: bring-up is about
ownership and ordering.

\section{Week 3 -- Raspberry Pi Radio Interface}

\paragraph{Done.} Ported the same register behaviour to Linux:
\texttt{/dev/spidev0.0} at 1~MHz, chip select, reset and DIO0 through
libgpiod \cite{libgpiod,spidevdocs}, DIO0 as a non-blocking event file
descriptor waited on with \texttt{poll()}, stale events drained before
transmit, a three-second transmit-complete wait with diagnostics,
retry on initialisation, and recovery when the chip leaves LoRa mode.

\paragraph{Problems.} The header gives DIO0 as GPIO~24 but the code
uses GPIO~23, and chip select (GPIO~22) is not in the header; the
kernel also toggles the unused CE0 pin; the Pi's rail browns out
under transmit current; \texttt{install.sh} installs neither libgpiod
nor SPI support.

\paragraph{Outcome.} Wiring follows the code; transmit power is at
the minimum and the radio is re-initialised when needed; a regulator
is the real fix. Lesson: a driver does not simply port to Linux.

\section{Week 4 -- Over-the-Air Integration}

\paragraph{Done.} Confirmed that the packet module is byte-identical
on both sides and that the register setups match apart from transmit
power; checked timing (500~ms slot timeout against a 225.5~ms longest
frame; six-second absence timeout against a 1.2~s idle cycle of 20
nodes); traced the acknowledgement flow.

\paragraph{Problems fixed.} The master stamped its own ID in POLL
packets, which nodes ignore; a 64-bit queue row ID cannot ride in an
8-bit field; a stale route echoed messages back to the sender.

\paragraph{Problems open (found by code review).} A node that sends
its own DATA or PRESENCE while a delivery rides on the POLL puts the
confirmation in \texttt{ack\_id}, which the master ignores for those
types, so the row stays queued and is delivered again. Redelivery
uses a new POLL ID, so the node's duplicate check cannot catch it:
delivery is at least once, not exactly once.

\section{Week 5 -- Phone Interface and Embedded Web Client}

\paragraph{Done.} Documented the captive portal (every DNS query
answered with 192.168.4.1; redirects for other hosts and paths), the
JSON WebSocket protocol, the offline-first client (IndexedDB with
contacts, chats and messages; pending queue flushed on reconnect;
back-off reconnect; three connection states; paging of 50 messages
and a 1\,000-message quota per chat), and the build step that inlines
CSS, JavaScript and base64 fonts into a 55\,642-byte gzip file
embedded in the firmware.

\paragraph{Problems.} The firmware still serves the small inline page,
not the embedded client; the client expects \texttt{ack} frames and
the firmware sends \texttt{queued} and \texttt{sent}; the chat list
inserts contact names and message previews with \texttt{innerHTML}
without escaping, a stored cross-site-scripting risk; the access
point accepts at most 10 phones.

\section{Week 6 -- Deployment, Operations and Field Validation}

\paragraph{Done.} Documented the systemd service, build and install
steps and the ESP-IDF toolchain; re-ran both host test suites on
26~November (node 43 of 43; master 149 checks, 0 failed); wrote a
staged field-validation plan (Chapter~6).

\paragraph{Problems.} Installation is not self-contained;
provisioning described in the README does not exist; credentials are
in the repository; configuration constants have drifted from the
radio registers; no measured results yet.

% ===================================================================
\chapter{Integration Findings}
% ===================================================================

\section{Errors, Causes and Status}

{\small
\begin{longtable}{|p{1.3in}|p{1.65in}|p{1.7in}|p{0.7in}|}
  \hline
  \textbf{Symptom} & \textbf{Cause} & \textbf{Fix or action} &
  \textbf{Status} \\
  \hline
  \endhead
  Chip stays in FSK mode &
  LoRa bit writable only in sleep &
  Sleep, set bit, standby; read back; retry & Fixed \\
  \hline
  Interrupt storms &
  Floating DIO0 &
  Pull-down; ignore empty IRQ; yield & Fixed \\
  \hline
  Nodes ignore POLL &
  Master stamped its own ID &
  POLL carries the addressed node ID & Fixed \\
  \hline
  Wrong delivery confirmed &
  64-bit row ID in 8-bit field &
  Pending record; delete on matching ACK & Fixed \\
  \hline
  Message echoed to sender &
  Stale route to sender's node &
  Detach route; store offline & Fixed \\
  \hline
  Chip leaves LoRa mode, TX hangs &
  Supply dip under transmit current &
  Minimum power; re-init; regulator advised & Mitigated \\
  \hline
  Board resets with WiFi on &
  Small regulator at full WiFi power &
  WiFi power limited to 8.5 dBm & Mitigated \\
  \hline
  DIO0 never fires on Pi &
  Header says GPIO 24, code uses 23 &
  Wire to the code; fix header & Open \\
  \hline
  Fresh Pi cannot build or find radio &
  \texttt{install.sh} lacks libgpiod and SPI &
  Add both to the script & Open \\
  \hline
  Logged RSSI too optimistic &
  Offset $-157$ used at 433 MHz &
  Use $-164$ before field tests & Open \\
  \hline
  Delivery repeated &
  Master ignores \texttt{ack\_id} on DATA and PRESENCE &
  Honour \texttt{ack\_id} on every reply & Open \\
  \hline
  Duplicates after lost ACK &
  Redelivery gets a new POLL ID &
  Stable end-to-end message ID & Open \\
  \hline
  Embedded UI not shown &
  Handler returns inline page &
  Serve embedded gzip file & Open \\
  \hline
  No delivered tick &
  Client expects \texttt{ack}, firmware sends \texttt{queued} &
  Align frames & Open \\
  \hline
  Markup in messages runs &
  Chat list uses \texttt{innerHTML} unescaped &
  Use \texttt{textContent} & Open \\
  \hline
  Cannot set node ID &
  No \texttt{/config} handler; ID compiled in &
  Add provisioning, store in NVS & Open \\
  \hline
  WiFi passphrase in repository &
  Defined in source and README &
  Change and externalise & Open \\
  \hline
  Config constants disagree &
  20 dBm in header, registers set 14 and 2 &
  Single source of truth & Open \\
  \hline
  Band and chip differ from design &
  SX1278 cannot reach 865 MHz &
  Decide band and hardware & Open \\
  \hline
\caption{Hardware, integration and interface findings}
\end{longtable}
}

\section{Patterns}

\begin{itemize}
  \item \textbf{Bugs live in meaning, not size.} A field of the right
    width with the wrong meaning passes unit tests on both sides and
    fails over the air.
  \item \textbf{Power is a software problem.} Three separate fixes
    trace back to supply dips under transmit current.
  \item \textbf{Documentation drifts.} Header pin, priority comments,
    README provisioning and configuration constants all disagree with
    the code.
\end{itemize}

% ===================================================================
\chapter{Phone Interface and Deployment}
% ===================================================================

\section{Captive Portal and Protocol}

A phone joining the WiFi network is sent to the chat page: a small DNS
server answers every query with the node address, and HTTP requests
for other hosts or paths are redirected to the root. The page opens a
WebSocket to \texttt{/ws}. The phone sends \texttt{attach} (phone
number) and \texttt{msg} (destination, text); the node answers
\texttt{attached}, \texttt{msg}, \texttt{queued}, \texttt{sent},
\texttt{status} (master online or not) and \texttt{error}. Sockets that
connect but never attach are closed after 30~s, and the HTTP server
purges the least recently used socket when full.

\section{Embedded Web Client}

The client is a 1\,271-line single-page application without
frameworks. It keeps identity locally, stores contacts, chats and
messages in IndexedDB (schema version 2), saves each outgoing message
as \emph{pending} before sending, flushes pending messages after
reconnect, reconnects with exponential back-off, and shows three
states (disconnected, node only, mesh live). The chat list is paged
(50 messages) with a 1\,000-message quota per chat. Its fonts are
embedded, so it needs no internet.

\section{Build Pipeline}

\texttt{build\_ui.py} inlines CSS and JavaScript into one HTML file,
converts the two fonts to base64 and compresses the result to a
55\,642-byte gzip file. CMake runs the script as a custom command
before compiling the component and embeds the file in the firmware
image \cite{espidf}.

\section{Access Point Settings}

One shared SSID on every node so that phones roam; WPA2 on channel 1;
at most 10 stations (a hardware limit); WiFi power limited to 8.5~dBm;
receive buffers, socket count, a 10~s task watchdog and the USB
console configured in \texttt{sdkconfig.defaults}.

\section{Master Deployment}

The master runs as a systemd service with automatic restart after five
seconds and logs to the journal through syslog; \texttt{make install}
and \texttt{install.sh} build, install and enable it. Node IDs are
command-line arguments, and the unit file starts the master with nodes
1 and 2. The queue database is an SQLite file in WAL mode
\cite{sqlitewal}.

% ===================================================================
\chapter{Status and Field-Validation Plan}
% ===================================================================

\section{Status}

\begin{itemize}
  \item Software layers complete and verified on a laptop: node suite
    43 of 43, master suite 149 checks with no failures (re-run on
    26 November).
  \item Hardware interface implemented on both platforms with
    diagnostics; no field measurements recorded yet.
\end{itemize}

\section{Validation Plan}

\begin{longtable}{|p{1.4in}|p{2.5in}|p{1.5in}|}
  \hline
  \textbf{Stage} & \textbf{Check} & \textbf{Expected (computed)} \\
  \hline
  \endhead
  1. Bring-up & Version register on both radios & \texttt{0x12} \\
  \hline
  2. Point to point & One frame each way at bench distance &
  Received, CRC good \\
  \hline
  3. Slot timing & POLL-to-ACK time & About 62 ms idle \\
  \hline
  4. Range & RSSI and SNR at rising distance, after fixing the
  RSSI offset & Margin about 34 dB up, 22 dB down at 1 km \\
  \hline
  5. Soak & Hours of polling with supply voltage logged &
  No brown-outs \\
  \hline
  6. Multi-node & Two nodes, roaming phones, offline delivery &
  Each message once \\
  \hline
\end{longtable}

\section{Fixes Before Deployment, in Priority Order}

\begin{enumerate}
  \item Honour \texttt{ack\_id} on every reply type at the master.
  \item Escape the chat-list fields in the web client.
  \item Serve the embedded client and align acknowledgement frames.
  \item Add node provisioning (ID in NVS).
  \item Complete \texttt{install.sh}; correct the DIO0 pin in
    \texttt{radio.h}.
  \item Correct the RSSI offset to $-164$.
  \item Change the WiFi passphrase and move it out of the repository.
  \item Settle the frequency band and radio hardware.
\end{enumerate}

\section{Reflection}

Passing software tests is the start of deployment work, not the end.
Radio band, power supply, antennas, installation scripts, credentials
and measurement each need their own checklist. The most useful habit
from this project was turning every claim into a number or a
checkable statement, because that is what exposed the mismatches
between design, code and documentation.

% ===================================================================
\printbibliography[heading=bibintoc]

\end{document}
